\chapter{Residue disks and the reduction of classes}
In this chapter $O$ is $\Zp$ or $\Oq = W(\Fq)$ unless stated otherwise; $K' :=
\mathrm{Frac}(O)$, with residue field $\kappa$. Recall the standing hypotheses:
$p \nmid \operatorname{disc} f$ and $p$ odd.

\begin{lemma}[Two kinds of points]
\label{lem:two-charts}
\lean{ChabautyColeman.point_int_or_inf}
\leanok
Every $P \in C(K')$ (affine or $\infty$) is exactly one of the following: of
\emph{finite type}, i.e.\ $P = (x_0, y_0)$ with $x_0, y_0 \in O$; or of \emph{infinite
type}, i.e.\ $v(x_0) = -2s < 0$, $v(y_0) = -(2g+1)s$, and
$t(P) := x_0^g/y_0 \in \mathfrak m$ (with $t(\infty) := 0$).
\end{lemma}
\begin{proof}
\leanok
If $v(x_0) < 0$ then $v(f(x_0)) = (2g+1)\,v(x_0)$ ($f$ is monic, so the leading term
dominates), hence $2v(y_0) = (2g+1)v(x_0)$ and
$2\,v(t(P)) = -v(x_0) > 0$, giving $t(P) \in \mathfrak m$ directly. If $v(x_0)\ge 0$
then $v(y_0^2) = v(f(x_0)) \ge 0$, so $y_0 \in O$.
\end{proof}

\begin{definition}[Reduction and residue disks]
\label{def:disk}
\lean{ChabautyColeman.redPoint, ChabautyColeman.ResidueDisk}
\leanok
\uses{lem:two-charts}
$\red(P) := (\bar x_0, \bar y_0) \in \bar C(\kappa)$ for points of finite type and
$\red(P) := \bar\infty$ for points of infinite type; the \emph{residue disk} of
$\bar P \in \bar C(\kappa)$ is $D_{\bar P} := \red^{-1}(\bar P) \subseteq C(K')$.
\end{definition}

\begin{lemma}[Square-root series]
\label{lem:sqrt-series}
\lean{ChabautyColeman.sqrtSeries}
\leanok
Let $c \in O[[z]]$ with $c(0)$ a square of a unit, $p$ odd. Then there is
$\eta \in O[[z]]$ with $\eta^2 = c$ and prescribed $\eta(0)$; explicitly, the binomial
series.
\end{lemma}
\begin{proof}
\leanok
The coefficient recursion
$\eta_n = \bigl(c_n - \sum_{0<i,j,\ i+j=n}\eta_i\eta_j\bigr)/(2\eta_0)$, solvable since
$2\eta_0 \in O^\times$, defines $\eta$; no completeness is needed. \emph{Formalization:}
\texttt{sqrtSeries} is constructed over an arbitrary commutative ring under the sole
hypothesis $2\eta_0 \in O^\times$.
\end{proof}

\begin{lemma}[Disk parametrizations]
\label{lem:disk-param}
\lean{ChabautyColeman.diskParamGeneric, ChabautyColeman.diskParamWeierstrass,
  ChabautyColeman.diskParamInf}
\leanok
\uses{def:disk, lem:sqrt-series, thm:hensel, lem:two-charts}
For each $\bar P \in \bar C(\kappa)$ there is a bijection
$z \colon D_{\bar P} \to \mathfrak m$ with inverse given by power series over $O$:
\begin{itemize}
\item[(i)] \emph{generic} $\bar P = (\bar a, \bar b)$ with $\bar b \ne 0$: fix the Hensel
lift $(a,b)$; then $z = x - a$ and $y = b\,\eta(z)$, with $\eta$ the square-root series
of $f(a+z)/f(a)$ (Lemma~\ref{lem:sqrt-series});
\item[(ii)] \emph{Weierstrass} $\bar P = (\bar\omega, 0)$: $z = y$ and
$x = \omega + \theta(z)$ with $\theta \in z^2 O[[z^2]]$ solving $z^2 = f(\omega +
\theta)$ by Theorem~\ref{thm:hensel} in series form; here $f'(\omega) \in O^\times$
because $\bar f$ is squarefree;
\item[(iii)] \emph{infinity}: $z = t = x^g/y$, and $x = z^{-2}\xi(z)$,
$y = z^{-(2g+1)}\psi(z)$ with $\xi, \psi \in 1 + zO[[z]]$ obtained by reverting
$t^2 = x^{2g}/f(x)$ in the variable $1/x$.
\end{itemize}
Weierstrass and generic disks over $\Fq$-points of the residue curve are treated over
$W(\Fq')$ for the field of definition $\Fq'$ of $\bar P$.
\end{lemma}
\begin{proof}
\leanok
In each chart, Hensel's lemma gives existence and uniqueness of the lift with prescribed
parameter value, and bijectivity onto $\mathfrak m$ follows from the uniqueness. For
(iii), revert by the standard inverse-series lemma in $O[[z]]$ (unit linear coefficient
after normalization). \emph{Formalization.} The Weierstrass and infinity charts need no
explicit power series in the parameter: for each parameter value the coordinate is the
unique Hensel root of a (in the $\infty$ case non-monic) polynomial equation.
\end{proof}

\begin{lemma}[Integral expansions of differentials]
\label{lem:omega-expansion}
\lean{ChabautyColeman.omegaExpansionGeneric, ChabautyColeman.omegaExpansionWeierstrass,
  ChabautyColeman.omegaExpansionInf}
\leanok
\uses{lem:disk-param}
Let $\omega = h(x)\,dx/y$ with $h \in O[x]$, $\deg h \le g-1$, and let $D$ be one of the
disks of Lemma~\ref{lem:disk-param}. Then $\omega|_{D} = w_D(z)\,dz$ with
$w_D \in O'[[z]]$ integral over the disk's base ring $O'$:
(i) $w = h(a+z)\,b^{-1}\eta(z)^{-1}$;
(ii) $w = 2h(x(z))/f'(x(z))$ times a unit series;
(iii) $w$ is the Laurent computation of
$h(z^{-2}\xi)\,d(z^{-2}\xi)\big/\bigl(z^{-(2g+1)}\psi\bigr)$, which is pole-free, of
order $2(g-1) - 2\deg h$ at lowest order.
\end{lemma}
\begin{proof}
\leanok
(i), (ii): chain rule (Lemma~\ref{lem:eval-hom}); the denominators are units because
$\bar f$ is squarefree and $p$ is odd. (iii): a valuation computation transcribing
$v_\infty(x^i\, dx/y) = 2g-2-2i \ge 0$; the leading coefficient is a unit multiple of
the leading coefficient of $h$. \emph{Formalization.} Each $w_D$ is a power series over
$O$ by construction, so integrality is definitional; the identity
``$\omega|_D = w_D\,dz$'' is a cleared-denominator differential identity in each chart,
and the order computation in (iii) is deferred to Lemma~\ref{lem:ord-match}.
\end{proof}

\begin{definition}[Disk antiderivatives]
\label{def:disk-antider}
\lean{ChabautyColeman.antideriv, ChabautyColeman.diskAntiderivGeneric,
  ChabautyColeman.diskAntiderivWeierstrass, ChabautyColeman.diskAntiderivInf}
\leanok
\uses{lem:omega-expansion}
$G^{\omega}_{D}(z) := \sum_k \tfrac{(w_D)_k}{k+1}\,z^{k+1}$ (so $G(0) = 0$), a
convergent function on $\mathfrak m$ by the estimates of
Theorem~\ref{thm:antideriv-zeros}.
\end{definition}

\begin{lemma}[Order matching]
\label{lem:ord-match}
\lean{ChabautyColeman.ordMatch_generic, ChabautyColeman.ordMatch_weierstrass,
  ChabautyColeman.ordMatch_inf,
  ChabautyColeman.isUnit_coeff_ordAffine_omegaExpansionGeneric,
  ChabautyColeman.isUnit_coeff_ordAffine_omegaExpansionWeierstrass,
  ChabautyColeman.isUnit_coeff_ordInf_omegaExpansionInf,
  ChabautyColeman.isUnit_coeff_of_order_map_eq}
\leanok
\uses{lem:omega-expansion, def:disk}
Let $\bar\omega := \bar h\,dx/y$ on $\bar C_{\Fpbar}$ with $\bar h \ne 0$. For every
$\bar P$:
$\ord_{\bar P}(\bar\omega) = \ord_z(\bar w_{D_{\bar P}}) =: m_{\bar P} < \infty$.
\end{lemma}
\begin{proof}
\leanok
The reduction of the expansion series is the expansion of the reduction, because the
parametrizations of Lemma~\ref{lem:disk-param} reduce to local parameters at $\bar P$ by
construction. The intrinsic order of the differential is computed on the Dedekind
coordinate ring $\A{\bar\kappa}$ at finite points --- the multiplicity of the point
ideal in $(\bar h)$, since $dx/y$ is regular and nonvanishing there --- and at $\infty$
as $\ord_\infty(\bar h) + (2g-2)$, with $\ord_\infty(\bar h) = -2\deg\bar h$ the order
at the infinite place (via the norm to $\bar\kappa(x)$) and $2g-2$ the contribution of
the canonical divisor. \emph{Formalization.} The form consumed by the per-disk count of
Theorem~\ref{thm:disk-bound} is the same statement read as a unit coefficient: in each
of the three charts, the integral expansion has a unit coefficient in degree
$m_{\bar P}$, because the reduced expansion has exactly that order.
\end{proof}

\begin{lemma}[Canonical degree over the residue field]
\label{lem:canon-deg}
\lean{ChabautyColeman.deg_div_omega_bar}
\leanok
\uses{def:hypring, lem:ord-match}
For $0 \ne \bar h$ of degree $\le g-1$ over $\Fpbar$:
$\deg \divisor(\bar h\,dx/y) = 2g-2$ on $\bar C_{\Fpbar}$.
\end{lemma}
\begin{proof}
\leanok
$\divisor(dx/y) = (2g-2)\,\bar\infty$, and $\divisor(\bar h)$ has degree $0$. The degree
is computed as the sum $\sum_{x_0}\sum_{y_0^2 = f(x_0)} \ord_{(x_0,y_0)}(\bar h)$ over
the affine points plus the order at $\infty$: each $x$-fibre contributes twice the root
multiplicity of $\bar h$ at $x_0$ (two points $(x_0,\pm b)$ generically, one doubled
Weierstrass point), so the affine part is $2\deg\bar h$, and the order at $\infty$ is
$2g-2-2\deg\bar h$.
\end{proof}

\begin{remark}[Formalization]
Lemma~\ref{lem:canon-deg} is formalized in
\texttt{HyperellipticCurves.Archive.CanonDegExact}; the main theorems do not depend on it,
since the count only needs the one-sided budget of Corollary~\ref{cor:sum-ord}.
\end{remark}

\begin{corollary}[Total vanishing budget]
\label{cor:sum-ord}
\lean{ChabautyColeman.sum_ord_le, ChabautyColeman.sum_ordPoint_le,
  ChabautyColeman.ordPoint, ChabautyColeman.finsum_ordPoint,
  ChabautyColeman.CurvePoint.equivOption, ChabautyColeman.ordAffine_eq_zero,
  ChabautyColeman.finite_support_ordPoint}
\leanok
\uses{lem:ord-match}
$\sum_{\bar P \in \bar C(\Fp)} m_{\bar P} \le 2g - 2$, with each $m_{\bar P} \ge 0$.
\end{corollary}
\begin{proof}
\leanok
The orders are nonnegative because $\bar\omega$ is regular
(Lemma~\ref{lem:omega-expansion} reduced), and the bound follows from a direct fibre
count over $\Fp$: each $x$-fibre has $0$, $1$ or $2$ rational points, hence contributes
at most twice the root multiplicity of $\bar h$ at $x_0$; summing the multiplicities
gives at most $\deg \bar h$; adding the order $2g - 2 - 2\deg\bar h$ at $\infty$ gives
$\le 2g-2$. The exact canonical degree (Lemma~\ref{lem:canon-deg}) is not needed.

\emph{Formalization.} The form summed in Theorem~\ref{thm:main-count} is: for \emph{any}
finite set $T$ of residue points, $\sum_{\bar P \in T} m_{\bar P} \le 2g-2$
(\texttt{sum\_ordPoint\_le}). The residue field is not assumed finite anywhere in this
chain: for $\bar h \neq 0$ the order vanishes at every point above a non-root of
$\bar h$, and each $x$-fibre is finite, so $\bar\omega$ has finitely many zeros and the
point sums are genuine finite sums over any field. This is what lets the counting
argument compose with Chapter~4, whose base ring has algebraically closed residue
field.
\end{proof}

\begin{corollary}[Vanishing budget of a space of differentials]
\label{cor:sum-ord-space}
\lean{ChabautyColeman.ordSpace, ChabautyColeman.sum_ordSpace_le,
  ChabautyColeman.exists_ordPoint_eq_ordSpace,
  ChabautyColeman.sum_ordPoint_le_two_mul_natDegree,
  ChabautyColeman.ordAffine_generic, ChabautyColeman.ordAffine_weierstrass,
  ChabautyColeman.ordAffine_le_ordAffine,
  Polynomial.finrank_add_natDegree_le_of_forall_dvd,
  Polynomial.exists_forall_natDegree_le_of_ne_bot,
  Polynomial.exists_dvd_forall_mem_and_rootMultiplicity_le}
\leanok
\uses{cor:sum-ord, lem:ord-match}
Let $V \ne 0$ be a subspace of the numerators of the regular differentials --- the
polynomials of degree $< g$ over $\Fpbar$ --- of codimension at most $r$, and put
\[
  n(V, \bar P) \;:=\; \min\{\ord_{\bar P}(\bar h\,dx/y) \;:\; 0 \ne \bar h \in V\} .
\]
Then $\sum_{\bar P \in T} n(V, \bar P) \le 2r$ for every finite set $T$ of points of
$\bar C$. For $\dim V = 1$ this is Corollary~\ref{cor:sum-ord}.
\end{corollary}
\begin{proof}
\leanok
For $x_0 \in \Fpbar$ let $\mu(x_0)$ be the smallest multiplicity of $x_0$ as a root of a
nonzero element of $V$, and let $q := \prod (x - x_0)^{\mu(x_0)}$, the product over the
$x$-coordinates of the affine points of $T$. Every nonzero $\bar h \in V$ is divisible by
$q$, since the factors are coprime and each divides $\bar h$. Let $D$ be the largest
degree of an element of $V$; then every element of $V$ lies in $q \cdot \{\deg \le D -
\deg q\}$, so $\dim V + \deg q \le D + 1$, and with $\dim V \ge g - r$,
\[ \deg q + \bigl(g - 1 - D\bigr) \;\le\; r . \]
At an affine point $\bar P$ over $x_0$ the order of $\bar h\,dx/y$ is the multiplicity of
$x_0$ in $\bar h$, doubled at a Weierstrass point, so it is monotone in that multiplicity
and $n(V, \bar P) = \ord_{\bar P}(q\,dx/y)$; the affine part of the sum is therefore at
most $2\deg q$ by the fibre count of Corollary~\ref{cor:sum-ord}, applied to $q$ with the
point at infinity left out. At infinity, the element of $V$ of maximal degree $D$ gives
$n(V, \infty) \le 2(g - 1 - D)$. Adding the two bounds gives $2r$.
\end{proof}

\begin{remark}[Formalization]
The minimum $n(V, \bar P)$ is \texttt{ordSpace}, an infimum over the nonzero elements of
$V$; it is attained (\texttt{exists\_ordPoint\_eq\_ordSpace}), which is how the per-disk
bound of Theorem~\ref{thm:disk-bound} is fed with an actual differential. The
displayed dimension inequality is \texttt{finrank\_add\_natDegree\_le\_of\_forall\_dvd},
stated for any subspace of polynomials of bounded degree with a common divisor: this is
the elementary substitute, available for a hyperelliptic curve, for the Riemann--Roch
argument of Theorem~6.4 of M.~Stoll, \emph{Independence of rational points on twists of
a given curve}, Compositio Math.\ 142 (2006), 1201--1214.
\end{remark}

\begin{lemma}[Content normalization of functions]
\label{lem:content}
\lean{ChabautyColeman.exists_primitive_smul, ChabautyColeman.exists_primitive_smul_pair,
  ChabautyColeman.exists_primitive_generator}
\leanok
\uses{lem:hypring-units}
Every $h \in \F{K'}^\times$ has a $K'^\times$-multiple $h_0 = a + by$ with
$a, b \in K'(x)$ admitting a \emph{primitive} presentation: numerator data in $O[x]$ of
unit content. The multiple is unique up to $O^\times$, since generators of a principal
ideal differ by $(\A{K'})^\times = K'^\times$.
\end{lemma}
\begin{proof}
\leanok
Scale by the content: primitivity is measured by the Gauss content
$\mathrm{cont}_v(p) = \sup_i v(p_i)$ on $K'[x]$; a common denominator $D$ clears
$h = a + by$ to $(P + Qy)/D$ with $P, Q, D \in K'[x]$, and scaling $D$ and the pair
$(P,Q)$ each to unit content --- attain the finite supremum and divide by the attaining
coefficient --- gives the primitive presentation. The scaling factor is a constant of
$K'^\times$. The statement holds for any valuation $v$; uniqueness up to $O^\times$ is
not needed downstream and is not formalized. \emph{Formalization.} Generators arrive in
the proof as explicit polynomial pairs $(P,Q)$, so only the pair-level normalizations
are used; the function-level statement is formalized in
\texttt{HyperellipticCurves.Archive.ContentGenerator}, and the main theorems do not depend on
it.
\end{proof}

\begin{theorem}[Specialization of principal divisors]
\label{thm:spec}
\leanok
\lean{ChabautyColeman.exists_center_finsum_emultiplicity_generic,
ChabautyColeman.exists_center_finsum_emultiplicity_weierstrass}
\uses{lem:content, lem:disk-param, thm:wprep, lem:two-charts}
Let $O_K$ be the ring of integers of a finite extension $K/K'$ and let
$h \in \F{K}^\times$ be primitive as in Lemma~\ref{lem:content}. Then $\bar h \ne 0$ in
$\F{\bar\kappa}$, and for every \emph{affine} $\bar P \in \bar C(\bar\kappa)$,
\[ \ord_{\bar P}(\bar h) \;=\; \sum_{P \in \divisor(h),\ \red(P) = \bar P} \ord_P(h), \]
the zeros of $h$ on the residue disk being enumerated as actual points over a splitting
discrete valuation ring. Affine points suffice: the reduction map of
Definition~\ref{def:redJ} takes values in the class group of the \emph{affine}
coordinate ring, where points escaping to $\bar\infty$ contribute the unit ideal, so the
statement at $\bar\infty$ is never needed.
\end{theorem}
\begin{proof}
\leanok
Write $h = P + Q\,y$ and expand it on the disk $D_{\bar P}$ (Lemma~\ref{lem:disk-param},
with charts over $O_K$). The number of zeros of the expansion in the open disk --- the
roots of the Weierstrass factor $W$ --- equals the vanishing order of its reduction
(Theorem~\ref{thm:wprep}); the per-disk point sums enumerate these zeros via Hensel
lifting of disk centres. At a Weierstrass point the hyperelliptic involution pairs the
disk parameters ($z \mapsto -z$), so the summed orders factor through the norm and no
individual conjugate order needs to be isolated. \emph{Formalization.} A chart-level
$\Cp$-version at every point including $\bar\infty$ is formalized in
\texttt{HyperellipticCurves.Archive.SpecializationGeneric} and
\texttt{HyperellipticCurves.Archive.SpecializationInf}; the main theorems do not depend on
it.
\end{proof}

\begin{definition}[Reduction of classes]
\label{def:redJ}
\lean{ChabautyColeman.redJ, ChabautyColeman.IsRedJPair, ChabautyColeman.exists_isRedJPair,
  ChabautyColeman.redIdealBar, ChabautyColeman.SplitSupported}
\leanok
\uses{lem:mumford-exists, def:disk, thm:spec}
Define $\redJ \colon \J{K'} \to \Cl(\A{\bar\kappa})$ as follows: represent
$c = \sum_i n_i\,\aj(P_i)$ with $P_i \in C(K_i)$ and $K_i/K'$ finite (possible by
Lemma~\ref{lem:mumford-exists}), and set $\redJ(c) := \sum_i n_i\,\aj(\red P_i)$.
\emph{Formalization.} The map is constructed over a complete discrete valuation ring
$O$ with algebraically closed residue field (the arithmetic version is the composite
with the base change to $\mathrm{Frac}\,O$): a class is represented over a finite
intermediate field $L$ of the algebraic closure by a \emph{split-supported} ideal ---
all prime divisors are maximal ideals of rational points; Mumford representatives become
split-supported over the splitting field of $u$ --- which is reduced pointwise to the
fixed residue curve, running the disk constructions over the integral closure $O_L$,
again a complete discrete valuation ring with the \emph{same} residue field.
\end{definition}

\begin{theorem}[Reduction is a well-defined homomorphism]
\label{thm:redJ-hom}
\lean{ChabautyColeman.redJHom, ChabautyColeman.isRedJPair_unique, ChabautyColeman.redJ_mul,
  ChabautyColeman.redJ_aj, ChabautyColeman.isPrincipal_redIdeal_span,
  ChabautyColeman.redIdealBar_map_of_splitSupported}
\leanok
\uses{def:redJ, thm:spec, lem:content}
$\redJ$ is well defined, is a group homomorphism, and satisfies
$\redJ(\aj(P)) = \aj(\red P)$.
\end{theorem}
\begin{proof}
\leanok
Independence of the representative: two representations differ by a principal divisor
$\divisor(h)$; normalize $h$ primitively (Lemma~\ref{lem:content}); by
Theorem~\ref{thm:spec} its pointwise reduction is $\divisor(\bar h)$, a principal
divisor, hence trivial in $\Cl(\A{\bar\kappa})$. Additivity is then formal.

\emph{Formalization.} Two split-supported representatives are compared over the
compositum and a further finite extension splitting the norms of the connecting
principal generators. There the reductions of the principal parts are principal ---
split the generator into a coprime primitive part, killed by Theorem~\ref{thm:spec},
and a split common factor, handled root by root through
$\mathfrak m_P\,\mathfrak m_{\iota P} = (x - \bar a)$ --- while the reductions of the
representatives are unchanged under base change: a split-supported ideal factors into
rational point ideals, which stay maximal and rational. Compatibility with the
Abel--Jacobi map: the reduction of the maximal ideal of an integral point is the maximal
ideal of the reduced point, and points reducing to $\bar\infty$ contribute the unit
ideal.
\end{proof}

\begin{lemma}[Frobenius-fixed classes are few]
\label{lem:frobfix-finite}
\lean{ChabautyColeman.frobFixed_finite, ChabautyColeman.HypJac.congr}
\leanok
\uses{lem:mumford-exists}
Let $\sigma$ be an automorphism of the base field with finitely many fixed elements that
fixes $\bar f$ (e.g.\ the $q$-power Frobenius on $\Fpbar$). Then the $\sigma$-fixed
subgroup of $\Cl(\A{\Fpbar})$ is finite: a fixed class has a $\sigma$-fixed
\emph{unique} reduced representative (Lemma~\ref{lem:mumford-exists} over $\Fpbar$),
hence Mumford data $(u,v)$ with coefficients in the fixed field, of which there are
finitely many.
\end{lemma}
\begin{proof}
\leanok
The automorphism induces an isomorphism of the coordinate ring and hence of the class
group (formalized in general: a ring isomorphism of Dedekind domains induces an
isomorphism of class groups, mapping the class of an ideal to the class of its image).
Uniqueness of the reduced Mumford representative transports fixedness of a class to
fixedness of its data; count monic $u$ of degree $\le g$ and $v$ of degree $< g$ with
coefficients in the fixed field.
\end{proof}

\begin{theorem}[The kernel of reduction has finite index]
\label{thm:J1-index}
\lean{ChabautyColeman.J1, ChabautyColeman.J1_finiteIndex, ChabautyColeman.redJ_fixed,
  ChabautyColeman.redIdealBar_congr, ChabautyColeman.adjoinRootSetEquiv,
  ChabautyColeman.congr_mapBase}
\leanok
\uses{thm:redJ-hom, lem:frobfix-finite}
$J^1(K') := \ker(\redJ)$ has finite index in $\J{K'}$, bounded by the number of Mumford
data $(u,v)$ over $\Fq$.
\end{theorem}
\begin{proof}
\leanok
Galois-equivariance of $\redJ$ on fixed classes: a class fixed by the arithmetic
automorphism has a fixed unique reduced Mumford representative $(u,v)$; its splitting
field $L = K(\text{roots of }u)$ is stable, and the reduction of the Mumford ideal over
$L$ transports along the induced automorphisms (a split-supported ideal factors into
rational point ideals, whose reductions are compatible with the semilinear
residue-field identification). Classes coming from $\J{K'}$ are fixed, since their
Mumford data live over $K'$, so $\J{K'}/J^1$ embeds into the group of Frobenius-fixed
classes, which is finite by Lemma~\ref{lem:frobfix-finite}.
\end{proof}

\begin{remark}[Formalization]
Since the pairing $\lambda$ of Definition~\ref{def:lambda} is defined on $J^1$
directly, the main theorems depend neither on Lemma~\ref{lem:frobfix-finite} nor on
Theorem~\ref{thm:J1-index}. Both are formalized, together with the equivariance apparatus
(a compatible family of automorphisms of $O$, its fraction field, the algebraic closure
and the residue field, all fixing $f$), in
\texttt{HyperellipticCurves.Archive.J1FiniteIndex},
\texttt{HyperellipticCurves.Archive.Equivariance} and
\texttt{HyperellipticCurves.Archive.StableSplitting}. The kernel $J^1$ itself, the base
change into it, and the isomorphism of class groups induced by a ring isomorphism are part
of the main development.
\end{remark}
